\documentclass[11pt]{article}
\PassOptionsToPackage{gray}{xcolor}
\usepackage[a4paper,margin=22mm]{geometry}
\usepackage{graphicx}
\usepackage{booktabs}
\usepackage{xcolor}
\usepackage{fontspec}
\usepackage[unicode,colorlinks=true,linkcolor=teal,urlcolor=teal]{hyperref}

\setmainfont[BoldFont=lmroman10-bold.otf]{lmroman10-regular.otf}
\definecolor{ReportInk}{HTML}{172B33}
\definecolor{ReportAccent}{HTML}{286C68}
\definecolor{ReportTint}{HTML}{EAF3F2}

\setlength{\parindent}{0pt}
\setlength{\parskip}{5pt}
\linespread{1.08}
\renewcommand{\arraystretch}{1.26}
\setlength{\tabcolsep}{5pt}

\newcommand{\reportlabel}[1]{{\large\color{ReportAccent}\bfseries #1}}
\newcommand{\reporttitle}[1]{%
  {\color{ReportInk}\Huge\bfseries #1\par}%
  \vspace{5mm}{\color{ReportAccent}\hrule height 1.3pt}%
}
\newcommand{\reportfigure}[3]{%
  \begin{center}
    {\bfseries #1}\par\smallskip
    \includegraphics[width=0.90\linewidth]{#2}\par\smallskip
    \begin{minipage}{0.9\linewidth}\small #3\end{minipage}
  \end{center}%
}

% Generated by python3 -m src.phase2_tables.
\newcommand{\EventN}{1237}
\newcommand{\PairedN}{843}
\newcommand{\NormalN}{847}
\newcommand{\OtherN}{390}
\newcommand{\UpperN}{622}
\newcommand{\LowerN}{615}
\newcommand{\MissingScoreN}{4}
\newcommand{\ThicknessCut}{2.319}
\newcommand{\TiedN}{5}
\newcommand{\CellAB}{488}
\newcommand{\CellAnB}{359}
\newcommand{\CellnAB}{134}
\newcommand{\CellnAnB}{256}
\newcommand{\ConditionalGap}{20.08}
\newcommand{\ExpectedJoint}{34.43}
\newcommand{\ExpectedCount}{425.9}
\newcommand{\JointGap}{5.02}
\newcommand{\BayesLikelihood}{57.62}
\newcommand{\BayesPosterior}{78.46}
\newcommand{\DrawN}{20}
\newcommand{\TailMin}{15}
\newcommand{\DrawFraction}{1.62}
\newcommand{\CountMean}{13.69}
\newcommand{\BinomialVar}{4.32}
\newcommand{\HypergeomVar}{4.25}
\newcommand{\BinomialTail}{35.95}
\newcommand{\HypergeomTail}{35.80}
\newcommand{\TailGap}{0.15}
\newcommand{\ScoreMin}{7}
\newcommand{\ScoreMax}{43}
\newcommand{\ScoreMean}{22.32}
\newcommand{\ScoreMedian}{22}
\newcommand{\ScoreFirstQuartile}{18}
\newcommand{\ScoreThirdQuartile}{26}
\newcommand{\ScoreEmpVar}{29.504}
\newcommand{\ScoreSampleVar}{29.539}
\newcommand{\ScoreSkew}{0.36}
\newcommand{\ScoreLowFlags}{0}
\newcommand{\ScoreHighFlags}{4}
\newcommand{\ScoreMode}{21}
\newcommand{\ScoreModeCount}{77}
\newcommand{\ScoreModePercent}{9.13}
\newcommand{\ThicknessMin}{1.474}
\newcommand{\ThicknessMax}{2.666}
\newcommand{\ThicknessMean}{2.336}
\newcommand{\ThicknessMedian}{2.338}
\newcommand{\ThicknessFirstQuartile}{2.2690}
\newcommand{\ThicknessThirdQuartile}{2.4175}
\newcommand{\ThicknessEmpVar}{0.014747}
\newcommand{\ThicknessSampleVar}{0.014764}
\newcommand{\ThicknessSkew}{-0.95}
\newcommand{\ThicknessLowFlags}{12}
\newcommand{\ThicknessHighFlags}{3}
\newcommand{\FirstPercentileObserved}{2.003}
\newcommand{\FirstPercentileNormal}{2.054}
\newcommand{\CovSample}{0.010466}
\newcommand{\CovEmpirical}{0.010453}
\newcommand{\DistinctPoints}{808}
\newcommand{\MargA}{68.47}
\newcommand{\MargnA}{31.53}
\newcommand{\MargB}{50.28}
\newcommand{\MargnB}{49.72}
\newcommand{\JointAB}{39.45}
\newcommand{\JointAnB}{29.02}
\newcommand{\JointnAB}{10.83}
\newcommand{\JointnAnB}{20.70}
\newcommand{\ConditionAB}{78.46}
\newcommand{\ConditionAnB}{58.37}

\graphicspath{{figures/phase2/}}
\hypersetup{
  pdftitle={Phase 2: Probability and distributions in the NACC teaching subset},
  allcolors=black,linkbordercolor={0 0 0},urlbordercolor={0 0 0}
}

\begin{document}
\reportlabel{Statistics course project, Track C (NACC)}
\vspace{7mm}

\reporttitle{Semantic verbal fluency and\\left inferior parietal cortical thickness}
\vspace{3mm}
{\Large Phase 2 | Probability and probability distributions\par}
\vspace{5mm}

{\large\bfseries Research question}\par
Among participants with normal cognition, how does the mean number of
correct responses in 60 seconds of animal naming differ per 0.1 mm greater
left inferior parietal cortical thickness, after accounting for age, sex,
and years of education? In this phase, we use probabilities to describe
the cognitive groups and examine the unadjusted relationship between
score and thickness among participants with normal cognition.

\section{Unit of observation and analysis samples}
We linked each person's UDS and MRI records using NACCID. Each file
contains one supplied record per participant. The Phase 1 build checks
for missing or repeated IDs before joining the files. The data do not
give visit dates or numbers, scan dates, or the time between measurements.
The sampling frame for the teaching subset is also undocumented.

\begin{center}\small
\begin{tabular}{@{}p{0.67\linewidth}r@{}}
\toprule
Analysis sample & Participants \\
\midrule
Usable UDS cognitive status and left cortical thickness, across all groups & \EventN{} \\
Normal cognition with a usable score and left cortical thickness & \PairedN{} \\
\bottomrule
\end{tabular}
\end{center}

Among the \NormalN{} participants with normal cognition, \MissingScoreN{}
had no usable animal-naming score. This leaves \PairedN{} people with
both measurements, whom we use for the two distributions and the covariance.

\section{Cross-modal events}
We define $A$ as normal cognition recorded in UDS (NACCUDSD = 1).
For the MRI event $B$, we group participants whose left inferior parietal
cortical thickness is at or above the median of all \EventN{} eligible
participants. This descriptive cutoff is \ThicknessCut{} mm. The
\TiedN{} measurements equal to the median are included in $B$.

\begin{center}\small
\begin{tabular}{@{}lrrr@{}}
\toprule
UDS status & At or above median & Below median & Total \\
\midrule
Normal cognition & \CellAB{} & \CellAnB{} & \NormalN{} \\
Other cognitive statuses & \CellnAB{} & \CellnAnB{} & \OtherN{} \\
\midrule
Total & \UpperN{} & \LowerN{} & \EventN{} \\
\bottomrule
\end{tabular}
\end{center}

\clearpage
\section{Empirical probabilities, independence, and Bayes' rule}
The table shows how often each event occurs, how often the events occur
together, and the share of participants with normal cognition in each
MRI group. Each empirical probability uses the count and denominator
shown alongside it.

\begin{center}\small
\begin{tabular}{@{}lrr@{}}
\toprule
Event & Count / reference group & Probability \\
\midrule
$A$ & \NormalN{} / \EventN{} & \MargA{}\% \\
$A^c$ & \OtherN{} / \EventN{} & \MargnA{}\% \\
$B$ & \UpperN{} / \EventN{} & \MargB{}\% \\
$B^c$ & \LowerN{} / \EventN{} & \MargnB{}\% \\
\midrule
$A\cap B$ & \CellAB{} / \EventN{} & \JointAB{}\% \\
$A\cap B^c$ & \CellAnB{} / \EventN{} & \JointAnB{}\% \\
$A^c\cap B$ & \CellnAB{} / \EventN{} & \JointnAB{}\% \\
$A^c\cap B^c$ & \CellnAnB{} / \EventN{} & \JointnAnB{}\% \\
\midrule
$A\mid B$ & \CellAB{} / \UpperN{} & \ConditionAB{}\% \\
$A\mid B^c$ & \CellAnB{} / \LowerN{} & \ConditionAnB{}\% \\
\bottomrule
\end{tabular}
\end{center}

The share with normal cognition was higher by \ConditionalGap{}
percentage points in the group at or above the thickness median.
We calculated the difference before rounding.

\subsection*{Empirical check of independence}
If the events were independent with these same marginal probabilities,
$\Pr(A)\Pr(B)$ would give a joint probability of \ExpectedJoint{}\%,
or \ExpectedCount{} people in the shared cell. We observed
\JointAB{}\% (\CellAB{} people). The observed share exceeds the
independence benchmark by \JointGap{} percentage points.

\subsection*{Updating the empirical probability}
Among people with normal cognition, \BayesLikelihood{}\% belong to $B$.
Using this information in Bayes' rule, knowing that a person belongs to
$B$ updates the probability of normal cognition from \MargA{}\% to
\BayesPosterior{}\%.
\[
\Pr(A\mid B)=\frac{\Pr(B\mid A)\Pr(A)}{\Pr(B)}
 =\frac{(\CellAB/\NormalN)(\NormalN/\EventN)}{\UpperN/\EventN}
 =\frac{\CellAB}{\UpperN}.
\]
This result matches the proportion obtained by counting directly in the table.

{\small References: Statistics Course Project Guide, Track C, pp. 8--10,
18--19. NACC teaching extract and variable definitions documented in the
Phase 1 project card.\par}

\clearpage
\section{A count model: binomial sampling}
Suppose we make \DrawN{} hypothetical draws from the \EventN{} eligible
records, giving each person an equal chance of selection. Let $Z$ be the
number of draws that select someone with normal cognition. With
replacement, the draws are independent and each has success probability
$p=\NormalN/\EventN$. The count then follows a binomial distribution,
with mean $\DrawN p$ and variance $\DrawN p(1-p)$. Without replacement,
each draw depends on who remains, and the count follows a hypergeometric
distribution.

\begin{center}\small
\begin{tabular}{@{}lrr@{}}
\toprule
Quantity & With replacement & Without replacement \\
\midrule
Expected count & \CountMean{} & \CountMean{} \\
Variance of the count & \BinomialVar{} & \HypergeomVar{} \\
$\Pr(Z\geq \TailMin)$ & \BinomialTail{}\% & \HypergeomTail{}\% \\
\bottomrule
\end{tabular}
\end{center}

We chose \TailMin{} as the example threshold, or three quarters of the
draws. Selecting \DrawN{} distinct participants covers \DrawFraction{}\%
of the reference set. The probability of reaching this threshold differs
by \TailGap{} percentage points between the two models. Both calculations
describe hypothetical sampling from this fixed set.

\section{Score and thickness distributions}
We summarize both measurements for the same \PairedN{} people.
The mean describes the center of the observed values. The empirical
variance is their average squared distance from that mean, using $n$
as the denominator. The sample variance uses $n-1$.

\begin{center}\small
\begin{tabular}{@{}lrr@{}}
\toprule
Statistic & Animal-naming score & Left thickness (mm) \\
\midrule
Range & \ScoreMin{}--\ScoreMax{} & \ThicknessMin{}--\ThicknessMax{} \\
Mean & \ScoreMean{} & \ThicknessMean{} \\
Median & \ScoreMedian{} & \ThicknessMedian{} \\
First and third quartiles & \ScoreFirstQuartile{}, \ScoreThirdQuartile{} &
\ThicknessFirstQuartile{}, \ThicknessThirdQuartile{} \\
Empirical variance ($n$) & \ScoreEmpVar{} & \ThicknessEmpVar{} \\
Sample variance ($n-1$) & \ScoreSampleVar{} & \ThicknessSampleVar{} \\
Empirical skewness & \ScoreSkew{} & \ThicknessSkew{} \\
Flagged below / above the fences & \ScoreLowFlags{} / \ScoreHighFlags{} &
\ThicknessLowFlags{} / \ThicknessHighFlags{} \\
\bottomrule
\end{tabular}
\end{center}

The variance units are squared responses for scores and mm squared for
thickness. We calculated quartiles by linear interpolation and flagged
values strictly outside the 1.5-IQR fences for review. These values are
included in all summaries.

\clearpage
\section{Empirical distributions}
\reportfigure{Figure 1. Empirical probabilities of animal-naming scores}{01_animal_score_pmf.pdf}{%
  Each bar shows the share of the \PairedN{} participants with that score.
  The most common score was \ScoreMode{}, recorded for \ScoreModeCount{}
  people (\ScoreModePercent{}\%). The distribution has a modest right tail.}

\reportfigure{Figure 2. Histogram of left inferior parietal cortical thickness}{02_thickness_histogram.pdf}{%
  Each bin spans 0.05 mm. Most measurements lie near the center, with a
  longer tail toward lower thickness values.}

\clearpage
\section{Normal reference and paired association}
\reportfigure{Figure 3. Normal Q-Q diagnostic for thickness}{03_thickness_normal_qq.pdf}{%
  We compare the ordered measurements with a normal reference fitted to
  the empirical mean and standard deviation, at positions $(i-0.5)/n$.
  The center follows the line more closely than the lower tail. The first
  percentiles are \FirstPercentileObserved{} mm in the data and
  \FirstPercentileNormal{} mm in the reference.}

\reportfigure{Figure 4. Thickness and animal-naming scores}{04_thickness_and_animal_scores.pdf}{%
  Each point represents one person. Some measurements overlap, so the
  \PairedN{} pairs occupy \DistinctPoints{} distinct coordinates. Scores
  vary widely even among people with similar thicknesses.}

The covariance is positive, at \CovSample{} mm times correct responses
with denominator $n-1$ and \CovEmpirical{} with denominator $n$.
It describes the unadjusted association in this observational subset.

\end{document}
